\section{Singapore deployment-scale techno-economic extension}
\label{sec:deployment-extension}
This bounded extension asks a different question from the original fixed-scale Gate-5 experiment. The original result remains 64 canonical nuclear-assisted cases at approximately 74.85 ktH$_2$/y with 0/64 joint CN4252 passes. Here the verified process is scaled as a complete industrial service---H$_2$ throughput, reforming, CCS, process heat, helium service, auxiliaries and annual cost---to test whether a larger literature-supported deployment can cross the \emph{absolute} 0.25 MtCO$_2$e/y threshold without exceeding S\$100/tCO$_2$e. No original Gate-5 range is altered.

\subsection{International nuclear and CCS cost basis}
Table~\ref{tab:deployment-evidence} separates source evidence from project conversions and scenario assumptions. JAEA's mature-design screen uses a 600 MWth GTHTR300-class unit and reports 0.7 JPY/MJ heat at 80\% availability, 40-year life and 3\% discount rate~\cite{jaea2014economics}; the GTHTR300C engineering precedent is 600 MWth at 950$^\circ$C with a 170 MW IHX for hydrogen cogeneration~\cite{jaeri2004gthtr300c}. Earlier GTHTR300 user requirements targeted 40--50 billion JPY/unit and greater than 90\% availability, but these are design targets rather than observed project costs~\cite{takei2006gthtr}. The independent modern bracket uses INL/GAIN's thermal-only HTGR-SMR meta-analysis~\cite{inlgain2024}. IEAGHG Case 1A supplies the incremental CCS capital anchor, 25-year life and 8\% discount basis~\cite{ieaghg2017smr}.

\begin{table}[htbp]
\centering\scriptsize
\caption{International nuclear/process-heat evidence used by the deployment-scale extension. JAEA values are mature-design/target estimates; INL values are modern meta-analysis screening points, not Singapore bids.}
\label{tab:deployment-evidence}
\resizebox{\linewidth}{!}{\pgfplotstabletypeset[col sep=comma,string type,
every head row/.style={before row=\toprule,after row=\midrule},every last row/.style={after row=\bottomrule}]
{../results/generated/deployment_evidence.csv}}
\end{table}

Project currency conversions are fixed screening assumptions (1 JPY=S\$0.0087, US\$1=S\$1.30, EUR1=S\$1.50). Incremental CCS capital is scaled linearly with captured-CO$_2$ throughput from the IEAGHG Case-1A capital increment; no unsupported economy-of-scale exponent is introduced. T\&S is S\$15/30/45 per tonne captured in the mature/central/adverse screens. The S\$15 low case is the project conversion of IEAGHG's EUR10/t source value; the higher cases are sensitivities. Singapore itself does not yet have an authoritative cross-border CCS tariff: MTI states that feasibility work is still obtaining clearer capture, transport and storage cost estimates~\cite{mticcs2025cost}. Singapore's 2026 S\$45/tCO$_2$e carbon tax is context only and is not subtracted from abatement cost~\cite{neacarbon2026}.

Because no authoritative Singapore nuclear-process integration cost exists, the extension uses explicit integration/site allowances of 10/20/30\% of allocated nuclear plus incremental CCS capital for Cases A/B/C. These are \textbf{SCREENING ASSUMPTIONS}, not predicted project costs.

\subsection{Required scale and reactor utilisation}
For a physical case with specific lifecycle abatement $\Delta e$, the annual threshold implies
\begin{equation}
M_{H_2,\min}=\frac{250{,}000\times1000}{\Delta e}.
\end{equation}
The best-abatement verified Gate-5 physical corner requires approximately the generated minimum in Table~\ref{tab:deployment-summary}; the conservative lifecycle corner has non-positive $\Delta e$ and therefore has no finite deployment scale that can satisfy the abatement threshold. The three economic scenarios are evaluated at 1.30 times the original \emph{annual} H$_2$ production, slightly above the minimum rather than exactly on the threshold.

\begin{table}[htbp]
\centering\scriptsize
\caption{Deployment-scale cases. Reactor MWth is installed module capacity; process heat is the service actually delivered to H$_2$ production.}
\label{tab:deployment-summary}
\resizebox{\linewidth}{!}{\pgfplotstabletypeset[col sep=comma,string type,
every head row/.style={before row=\toprule,after row=\midrule},every last row/.style={after row=\bottomrule}]
{../results/generated/deployment_summary.csv}}
\end{table}

Figure~\ref{fig:deployment-abatement} shows why scale matters for the absolute CN4252 criterion. It does not improve specific lifecycle performance; it only multiplies a positive specific abatement over a larger useful industrial service.

\begin{figure}[htbp]\centering
\begin{tikzpicture}
\begin{axis}[width=.88\linewidth,height=.50\linewidth,xlabel={Annual H$_2$ production (t/y)},ylabel={Annual avoided emissions (tCO$_2$e/y)},grid=major,scaled x ticks=false,xticklabel style={/pgf/number format/fixed,/pgf/number format/1000 sep={,}}]
\addplot+[thick,mark=*] table[col sep=comma,x=h2_t_y,y=annual_avoided_t]{../results/generated/deployment_scale_curve.csv};
\draw[dashed,thick] (axis cs:70000,250000)--(axis cs:225000,250000);
\end{axis}
\end{tikzpicture}
\caption{Deployment scale versus annual lifecycle abatement for the verified positive-abatement physical corner. The dashed line is the 0.25 MtCO$_2$e/y threshold. The conservative non-positive-abatement corner is excluded because scaling a negative specific abatement cannot create a positive annual benefit.}
\label{fig:deployment-abatement}
\end{figure}

\subsection{Forward economics and allocation}
Annual candidate and baseline costs are constructed before applying the S\$100/t criterion. Hydrogen-only cases pay the full 600 MWth module service/capital. Cogeneration cases allocate reactor capital/O\&M in proportion to the process-heat thermal-capacity share; the balance of reactor output must therefore have a useful electricity/steam/heat customer and is not assigned zero cost. This is a capacity-allocation screen, not a revenue forecast.

\begin{figure}[htbp]\centering
\begin{tikzpicture}
\begin{axis}[width=.9\linewidth,height=.52\linewidth,xlabel={Annual scale relative to original plant},ylabel={Forward abatement cost (S\$/tCO$_2$e)},grid=major,ymin=0,ymax=500]
\addplot+[thick] table[col sep=comma,x=annual_scale,y=cost_sgd_t,restrict expr to domain={\thisrow{scenario_id}}{0:0}]{../results/generated/deployment_cost_curve.csv};\addlegendentry{JAEA mature}
\addplot+[thick] table[col sep=comma,x=annual_scale,y=cost_sgd_t,restrict expr to domain={\thisrow{scenario_id}}{1:1}]{../results/generated/deployment_cost_curve.csv};\addlegendentry{Modern central}
\addplot+[thick] table[col sep=comma,x=annual_scale,y=cost_sgd_t,restrict expr to domain={\thisrow{scenario_id}}{2:2}]{../results/generated/deployment_cost_curve.csv};\addlegendentry{FOAK/adverse}
\draw[dashed,thick] (axis cs:1,100)--(axis cs:3,100);
\end{axis}
\end{tikzpicture}
\caption{Deployment scale versus forward abatement cost for cogeneration allocation. The horizontal line is S\$100/tCO$_2$e. Scaling alone does not rescue an expensive cost regime: the mature JAEA case is below the threshold, whereas the modern central and adverse screens remain above it.}
\label{fig:deployment-cost}
\end{figure}

Figure~\ref{fig:deployment-feasibility} presents the joint feasibility envelope in terms of annual scale and allowable delivered nuclear-heat cost for the low T\&S/mature integration structure. The vertical condition is set by annual abatement; the curve gives the heat price that would make total forward cost exactly S\$100/tCO$_2$e. The JAEA converted heat value is shown as literature evidence, not as a Singapore observed price.

\begin{figure}[htbp]\centering
\begin{tikzpicture}
\begin{axis}[width=.88\linewidth,height=.52\linewidth,xlabel={Annual scale relative to original plant},ylabel={Delivered nuclear heat (S\$/GJ)},grid=major,ymin=0,ymax=30]
\addplot+[thick] table[col sep=comma,x=annual_scale,y=break_even_heat_sgd_gj]{../results/generated/deployment_heat_feasibility.csv};\addlegendentry{S\$100/t break-even}
\draw[dashed,thick] (axis cs:1.284,0)--(axis cs:1.284,30);
\draw[densely dashed,thick] (axis cs:1,6.09)--(axis cs:3,6.09);
\addlegendimage{densely dashed,thick}\addlegendentry{JAEA 0.7 JPY/MJ conversion}
\end{axis}
\end{tikzpicture}
\caption{Joint-feasibility envelope for the mature low-T\&S structure. A joint pass requires scale to the right of the abatement boundary and heat cost below the break-even curve. This is a deterministic feasibility boundary, not a probability map.}
\label{fig:deployment-feasibility}
\end{figure}

\subsection{Principal deployment result}
Table~\ref{tab:deployment-principal} gives the three literature/scenario classes at the 1.30 annual scale. The original 0/64 result is unchanged because these are new deployment-scale cases.

\begin{table}[htbp]\centering\scriptsize
\caption{Principal Singapore deployment-scale results for cogeneration allocation.}
\label{tab:deployment-principal}
\resizebox{\linewidth}{!}{\pgfplotstabletypeset[col sep=comma,string type,
every head row/.style={before row=\toprule,after row=\midrule},every last row/.style={after row=\bottomrule}]
{../results/generated/deployment_principal.csv}}
\end{table}

\begin{figure}[htbp]\centering
\begin{tikzpicture}
\begin{axis}[ybar,width=.9\linewidth,height=.5\linewidth,ylabel={Annual incremental cost (S\$)},symbolic x coords={NG and auxiliary net,Nuclear heat service,CCS capital annualisation,Integration/site annualisation,CO2 transport-storage},xtick=data,x tick label style={rotate=30,anchor=east,font=\scriptsize},grid=major,scaled y ticks=false,yticklabel style={/pgf/number format/fixed,/pgf/number format/1000 sep={,}}]
\addplot table[col sep=comma,x=component,y=annual_sgd]{../results/generated/deployment_cost_breakdown.csv};
\end{axis}
\end{tikzpicture}
\caption{Annual incremental-cost breakdown for the conditional JAEA mature cogeneration case. The negative NG/auxiliary-net bar is an avoided-cost contribution, not a subsidy. Reactor heat, CCS capital, integration allowance and T\&S are charged independently to avoid double counting.}
\label{fig:deployment-cost-breakdown}
\end{figure}

The extension therefore identifies a \textbf{CONDITIONAL SCENARIO PASS}: at the larger 1.30 annual H$_2$ scale, the JAEA mature-design cogeneration allocation passes both thresholds under the declared low T\&S and integration assumptions. The same physical deployment fails the cost criterion under the modern-central and FOAK/adverse nuclear-cost regimes, and the JAEA hydrogen-only allocation also fails because unused reactor capacity cannot be treated as free. This does not demonstrate commercial feasibility. It shows that a passing region exists only when mature-design nuclear heat economics, useful cogeneration of the remaining reactor output and low-end CCS service costs are simultaneously realised.
